% from ega4-2-fr.tex
\oldpage[IV-2]{203}

tels que $f$ soit fidèlement plat; comme par hypothèse $g\circ f$ est un
$\boldsymbol P$-morphisme, il en est de même de $g$ en vertu de
$(\mathrm P_{\mathrm{II}})$.
\end{proof}

\begin{corollary}[7.4.7]
\label{IV.2.s7.4.7-fr}
Supposons que $\boldsymbol P$ vérifie les conditions
$(\mathrm P_{\mathrm I})$, $(\mathrm P'_{\mathrm I})$,
$(\mathrm P_{\mathrm{II}})$, $(\mathrm P_{\mathrm{III}})$ et
$(\mathrm P_{\mathrm{IV}})$. Alors, si $A$ est un
$\boldsymbol P$-anneau (7.4.5), le morphisme canonique
$\operatorname{Spec}(A[[T_1,\ldots,T_r]])\to\operatorname{Spec}(A)$ est un
$\boldsymbol P$-morphisme. En particulier, si de plus $A$ est intègre, $K$
son corps des fractions, et si $\mathfrak p$ est un idéal premier de
$B=A[[T_1,\ldots,T_r]]$ tel que $\mathfrak p\cap A=0$, alors la propriété
$\boldsymbol P(B_{\mathfrak p},K)$ est vraie.
\end{corollary}

\begin{proof}
En effet, le morphisme canonique
$\operatorname{Spec}(B)\to\operatorname{Spec}(A)$ se factorise en
\[
\operatorname{Spec}(A[[T_1,\ldots,T_r]])
\xrightarrow{f}\operatorname{Spec}(A[T_1,\ldots,T_r])
\xrightarrow{g}\operatorname{Spec}(A).
\]
Il est clair que le morphisme $g$ est régulier (0, 17.3.7); en vertu de
(7.4.5), $A[T_1,\ldots,T_r]$ est un $\boldsymbol P$-anneau, donc il résulte
de (7.4.6) que $f$ est un $\boldsymbol P$-morphisme; comme $g$ est régulier,
donc est aussi un $\boldsymbol P$-morphisme (7.3.5, (i)), il en est de même
de $g\circ f$ en vertu de $(\mathrm P'_{\mathrm I})$.

On notera que la conclusion est encore valable si au lieu de supposer que
$\boldsymbol P$ vérifie $(\mathrm P'_{\mathrm I})$ et que tout morphisme
régulier est un $\boldsymbol P$-morphisme, on suppose seulement qu'un
morphisme composé $g\circ f$ est un $\boldsymbol P$-morphisme lorsque $g$
est régulier et $f$ un $\boldsymbol P$-morphisme (condition symétrique en
quelque sorte de $(\mathrm P'_{\mathrm I})$).
\end{proof}

\begin{remark}[7.4.8]
\label{IV.2.s7.4.8-fr}
Les résultats précédents posent les problèmes suivants :
\begin{enumerate}
\item[\emph{A)}] Soit $A$ un anneau de Zariski \emph{complet}, et soit
$\mathfrak J$ un idéal de définition de $A$; si l'anneau $A/\mathfrak J$
est un $\boldsymbol P$-anneau, en est-il de même de $A$? Il en résulterait
que pour tout $\boldsymbol P$-anneau noethérien $A$ et tout idéal
$\mathfrak J$ de $A$, le séparé complété $\widehat A$ pour la topologie
$\mathfrak J$-préadique serait aussi un $\boldsymbol P$-anneau.
\item[\emph{B)}] Soit $k$ un corps valué complet non discret; on appelle
encore anneau des \emph{séries formelles restreintes}
$k\{T_1,\ldots,T_n\}$ le sous-anneau de l'anneau de séries formelles
$k[[T_1,\ldots,T_n]]$ formées des séries dont les coefficients \emph{tendent
vers 0}. Un tel anneau est-il un $\boldsymbol P$-anneau?
\item[\emph{C)}] Si $A$ est un $\boldsymbol P$-anneau linéairement
topologisé, $S$ une partie multiplicative de $A$, les anneaux
$A\{S^{-1}\}$ (0$_{\mathrm I}$, 7.6.1) et $A_{\{S\}}$
(0$_{\mathrm I}$, 7.6.15) sont-ils des $\boldsymbol P$-anneaux?
\end{enumerate}
\end{remark}

\subsection{\texorpdfstring{Un critère pour les $\boldsymbol P$-morphismes}{Un critère pour les P-morphismes}}
\label{subsection:IV.7.5-fr}

\begin{env}[7.5.0]
\label{IV.2.s7.5.0-fr}
Ce numéro ne sera pas utilisé dans la suite du chapitre IV, et peut donc être
omis en première lecture. Nous verrons d'ailleurs plus loin (7.9.8) que les
résultats du présent numéro peuvent être considérablement améliorés quand on
dispose de la « résolution des singularités ».

Dans la suite de ce numéro, nous considérerons une propriété
$\boldsymbol R(A)$, et nous désignerons par $\boldsymbol P(Z,k)$ la propriété
suivante :
\begin{quote}
« $Z$ est un préschéma localement noethérien sur un corps $k$, et, pour toute
extension finie $k'$ de $k$, si l'on pose $Z'=Z\otimes_k k'$, la propriété
$\boldsymbol R(\mathscr O_{Z',z'})$ est vraie pour tout $z'\in Z'$. »
\end{quote}
\end{env}

\begin{theorem}[7.5.1]
\label{IV.2.s7.5.1-fr}
Soient $A$, $B$ deux anneaux locaux noethériens complets, $\mathfrak m$
l'idéal maximal de $A$, $k=A/\mathfrak m$ son corps résiduel,
$\varphi:A\to B$ un homomorphisme local tel que :
\begin{enumerate}
\item[\emph{(i)}] Le corps résiduel de $B$ est une extension finie de $k$.
\item[\emph{(ii)}] $B$ est un $A$-module plat.
\end{enumerate}
Soit d'autre part $\boldsymbol R(C)$ une propriété vérifiant la condition
$(\mathrm R_{\mathrm{III}})$ (7.3.18) et la condition suivante :
\begin{quote}
$(\mathrm R_{\mathrm{IV}})$ Pour tout anneau local $C$ en un idéal premier
d'un anneau local noethérien complet et tout élément $C$-régulier $t$ dans
l'idéal maximal de $C$, la propriété $\boldsymbol R(C/tC)$ implique
$\boldsymbol R(C)$.
\end{quote}

